
\chapter*{参考文献}
\phantomsection\addcontentsline{toc}{chapter}{参考文献}
\markboth{参考文献}{参考文献}
\begingroup
\small\raggedright
\interlinepenalty=10000
\setlength{\emergencystretch}{2em}
\begin{list}{}{\leftmargin=2.1em \labelwidth=0pt \labelsep=0pt
  \itemindent=-2.1em \parsep=0pt \itemsep=2pt \topsep=0pt}
\item Ajami, S. and Teimouri, F. (2015). Features and application of wearable biosensors in medical care. Journal of Research in Medical Sciences: The Official Journal of Isfahan University of Medical Sciences 20:1208--1215.
\item Akutekwe, A. and Seker, H. (2015). Inference of nonlinear gene regulatory networks through optimized ensemble of support vector regression and dynamic Bayesian networks. Engineering in Medicine and Biology Society (EMBC), 2015 37th Annual International Conference of the IEEE 8177--8180.
\item Albarqouni, S., Baur, C., Achilles, F., Belagiannis, V., Demirci, S. and Navab, N. (2016). Aggnet: Deep learning from crowds for mitosis detection in breast cancer histology images. IEEE Transactions on Medical Imaging 35:1313--1321.
\item Alonso, A. M., Casado, D. and Romo, J. (2012). Supervised classification for functional data: A weighted distance approach. Computational Statistics \& Data Analysis 56: 2334--2346.
\item Anders, S. and Huber, W. (2010). Differential expression analysis for sequence count data. Genome Biol 11:R106.
\item Anderson, T. W. (1984). An introduction to multivariate statistical analysis. 2nd ed., John Wiley \& Sons, New York.
\item Andrew, N. (2011). Sparse autoencoder. CS294A Lecture Notes 72.
\item Aschard, H., Vilhjálmsson, B. J., Greliche, N., Morange, P.-E., Trégouët, D.-A. and Kraft, P. (2014). Maximizing the power of principal-component analysis of correlated phenotypes in genome-wide association studies. The American Journal of Human Genetics 94:662--676.
\item Bühlmann, P., Peters, J. and Ernest, J. (2014). CAM: Causal additive models, high-dimensional order search and penalized regression. The Annals of Statistics 42: 2526--2556.
\item Bailey, D. L., Townsend, D. W., Valk, P. E. and Maisey, M. N. (2005). Positron emission tomography. Springer.
\item Baiocchi, M., Cheng, J. and Small, D. S. (2014). Instrumental variable methods for causal inference. Statistics in Medicine 33:2297--2340.
\item Balevich, E. C., Haznedar, M. M., Wang, E. et al. (2015). Corpus callosum size and diffusion tensor anisotropy in adolescents and adults with schizophrenia. Psychiatry Research: Neuroimaging 231:244--251.
\item Bartlett, M. and Cussens, J. (2013). Advances in Bayesian network learning using integer programming. arXiv:1309.6825.
\item Bartlett, M. S. (1939). A note on tests of significance in multivariate analysis. Mathematical Proceedings of the Cambridge Philosophical Society 180--185.
\item Barzel, B. and Barabási, A.-L. (2013). Network link prediction by global silencing of indirect correlations. Nature Biotechnology 31:720--725.
\item Belilovsky, E., Varoquaux, G. and Blaschko, M. B. (2016). Testing for differences in Gaussian graphical models: Applications to brain connectivity. arXiv:1512.08643.
\item Besag, J. (1977). Efficiency of pseudolikelihood estimation for simple Gaussian fields. Biometrika 616--618.
\item Bollen, K. A. (1989). Structural equations with latent variables. John Wiley \& Sons, New York.
\item Bollen, K. A. (2012). Instrumental variables in sociology and the social sciences. Annual Review of Sociology 38:37--72.
\item Bolstad, B. M., Irizarry, R. A., Åstrand, M. and Speed, T. P. (2003). A comparison of normalization methods for high density oligonucleotide array data based on variance and bias. Bioinformatics 19:185--193.
\item Bouvrie, J. (2006). Notes on convolutional neural networks. \url{http://cogprints.org} /5869/1/cnn\_tutorial.pdf.
\item Boyd, S., Parikh, N., Chu, E., Peleato, B. and Eckstein, J. (2011). Distributed optimization and statistical learning via the alternating direction method of multipliers. Foundations and Trends® in Machine Learning 3:1--122.
\item Boyle, E. A., Li, Y. I. and Pritchard, J. K. (2017). An Expanded View of Complex Traits: From Polygenic to Omnigenic. Cell 169:1177--1186.
\item Brauns, S., Gollub, R. L., Walton, E. et al. (2013). Genetic variation in GAD1 is associated with cortical thickness in the parahippocampal gyrus. Journal of Psychiatric Research 47:872--879.
\item Bressler, S. L. and Seth, A. K. (2011). Wiener--Granger causality: A well established methodology. Neuroimage 58:323--329.
\item Brito, C. and Pearl, J. (2002). Generalized instrumental variables. Proceedings of the Eighteenth Conference on Uncertainty in Artificial Intelligence 85--93.
\item Cai, X., Bazerque, J. A. and Giannakis, G. B. (2013). Inference of gene regulatory networks with sparse structural equation models exploiting genetic perturbations. PLoS Comput Biol 9:e1003068.
\item Callaway, E. (2017). Genome studies attract criticism. NATURE 546:463.
\item Cancer Genome Atlas Research Network. (2011). Integrated genomic analyses of ovarian carcinoma. NATURE 474:609--615.
\item Candès, E. J., Li, X., Ma, Y. and Wright, J. (2011). Robust principal component analysis? Journal of the ACM (JACM) 58:11.
\item Cascio, D., Magro, R., Fauci, F., Iacomi, M. and Raso, G. (2012). Automatic detection of lung nodules in CT datasets based on stable 3D mass--spring models. Computers in Biology and Medicine 42:1098--1109.
\item Chang, C. C. and Lin, C. J. (2011). LIBSVM: A library for support vector machines. ACM Transactions on Intelligent Systems and Technology 2:1--27.
\item Chapter 2: Variational Bayesian theory. www.cse.buffalo.edu/faculty/mbeal/thesis /beal03\_2.pdf.
\item Chen, B. and Pearl, J. (2014). Graphical tools for linear structural equation modeling. UCLA Cognitive Systems Laboratory, Technical Report (R-432).
\item Chen, X., Zou, C. and Cook, R. D. (2010). Coordinate-independent sparse sufficient dimension reduction and variable selection. The Annals of Statistics 3696--3723.
\item Chen, Z., Zhang, K. and Chan, L. (2013). Nonlinear causal discovery for high dimensional data: A kernelized trace method. Data Mining (ICDM), 2013 IEEE 13th International Conference on 1003--1008.
\item Cheng, J., Grainer, G., Kelly, J., Bell, D. and Lius, W. (2002). Learning bayesian networks from data: An information-theory based approach. Artificial Intelligence 137:43--90.
\item Cheng, S., Guo, M., Wang, C., Liu, X., Liu, Y. and Wu, X. (2016). MiRTDL: A deep learning approach for miRNA target prediction. IEEE/ACM Transactions on Computational Biology and Bioinformatics 13:1161--1169.
\item Ciresan, D., Giusti, A., Gambardella, L. M. and Schmidhuber, J. (2012). Deep neural networks segment neuronal membranes in electron microscopy images. Advances in Neural Information Processing Systems 2843--2851.
\item Clemmensen, L., Hastie, T., Witten, D. and Ersbøll, B. (2011). Sparse discriminant analysis. Technometrics 53:406--413.
\item Clyde, D. (2017). Disease genomics: Transitioning from association to causation with eQTLs. Nature Reviews Genetics 18:271.
\item Cover, T. M. and Thomas, J. A. (2012). Elements of information theory. John Wiley \& Sons, New York.
\item Cowley, B. and Vinci, G. (2014). Summary and discussion of: ``Brownian Distance Covariance''. \url{http://www.stat.cmu.edu/~ryantibs/journalclub/dcov.pdf}.
\item Cussens, J. (2012). Bayesian network learning with cutting planes. arXiv:1202.3713.
\item Cussens, J. (2014). Integer programming for Bayesian network structure learning. Quality Technology \& Quantitative Management 11:99--110.
\item Danaher, P., Wang, P. and Witten, D. M. (2014). The joint graphical lasso for inverse covariance estimation across multiple classes. Journal of the Royal Statistical Society: Series B (Statistical Methodology) 76:373--397.
\item Daniusis, P., Janzing, D., Mooij, J. et al. (2012). Inferring deterministic causal relations. Proceedings of the 26th Conference on Uncertainty in Artificial Intelligence (UAI) pp. 1--8.
\item Das, K., Li, J., Wang, Z. et al. (2011). A dynamic model for genome-wide association studies. Human genetics 129:629--639.
\item De Campos, L. M. and Huete, J. F. (2000). A new approach for learning belief networks using independence criteria. International Journal of Approximate Reasoning 24:11-- 37.
\item Delude, C. M. (2015). Deep phenotyping: The details of disease. NATURE 527:S14--S15.
\item Dhawan, A. P. (2011). Medical image analysis. Vol. 31. John Wiley \& Sons.
\item Dippel, C., Gold, R., Heblich, S. and Pinto, R. (2017). Instrumental Variables and Causal Mechanisms: Unpacking The Effect of Trade on Workers and Voters. National Bureau of Economic Research.
\item Dutta, M. (2015). Assessment of feature extraction techniques for hyperspectral image classification. Computer Engineering and Applications (ICACEA), 2015 International Conference on Advances in 499--502.
\item Ernest, J. (2016). Causal inference in semiparametric and nonparametric structural equation models, Ph.D. Thesis, ETH Zurich.
\item Fan, X., Malone, B. and Yuan, C. (2014). Finding optimal Bayesian network structures with constraints learned from data. Proceedings of the Thirtieth Conference on Uncertainty in Artificial Intelligence 200--209.
\item Farabet, C., Couprie, C., Najman, L. and LeCun, Y. (2013). Learning hierarchical features for scene labeling. IEEE Transactions on Pattern Analysis and Machine Intelligence 35:1915--1929.
\item Fortin, J. P., Triche, T. J., Jr. and Hansen, K. D. (2017). Preprocessing, normalization and integration of the Illumina Human Methylation EPIC array with minfi. Bioinformatics 33:558--560.
\item Friedman, J., Hastie, T. and Tibshirani, R. (2008). Sparse inverse covariance estimation with the graphical lasso. Biostatistics 9:432--441.
\item Friedman, N. and Nachman, I. 2000. Gaussian process networks. In Proceedings of the Sixteenth Conference on Uncertainty in Artificial Intelligence. Stanford, CA: Morgan Kaufmann Publishers Inc.
\item Friston, K. J., Harrison, L. and Penny, W. (2003). Dynamic causal modelling. Neuroimage 19:1273--1302.
\item Fügenschuh, A. and Martin, A. (2005). Computational integer programming and cutting planes. Handbooks in Operations Research and Management Science 12:69-- 121.
\item Fusi, N. and Listgarten, J. (2016). Flexible modelling of genetic effects on function-valued traits. Journal of Computational Biology 24:524--535.
\item Gamboa, J. C. B. (2017). Deep learning for time-series analysis. arXiv:1701.01887.
\item Gao, W., Emaminejad, S., Nyein, H. Y. Y. et al. (2016). Fully integrated wearable sensor arrays for multiplexed in situ perspiration analysis. NATURE 529:509--514.
\item Garcia-Garcia, A., Orts-Escolano, S., Oprea, S., Villena-Martinez, V. and Garcia-Rodriguez, J. (2017). A review on deep learning techniques applied to semantic segmentation. arXiv:1704.06857.
\item Garin-Muga, A. and Borro, D. (2014). Review and challenges of brain analysis through DTI measurements. Studies in Health Technology and Informatics 207:27--36.
\item Gene expression. \url{https://www2.stat.duke.edu/courses/Spring04/sta278/refinfo/Gene} \_Expression.pdf.
\item Gershman, S. J. and Daw, N. D. (2017). Reinforcement learning and episodic memory in humans and animals: An integrative framework. Annual Review of Psychology 68:101--128.
\item Gianola, D. and Sorensen, D. (2004). Quantitative genetic models for describing simultaneous and recursive relationships between phenotypes. Genetics 167:1407--1424.
\item Gibbs, A. L. and Su, F. E. (2002). On choosing and bounding probability metrics. International Statistical Review 70:419--435.
\item Gibson, E., Hu, Y., Huisman, H. J. and Barratt, D. C. (2017). Designing image segmentation studies: Statistical power, sample size and reference standard quality. Med Image Anal 42:44--59.
\item Glover, G. H. (2011). Overview of functional magnetic resonance imaging. Neurosurgery Clinics of North America 22:133--139.
\item González, I. (2014). Tutorial: Statistical analysis of RNA-Seq data. www.nathalievilla .org/doc/pdf/tutorial-rnaseq.pdf.
\item Granger, C. W. (1969). Investigating causal relations by econometric models and cross-spectral methods. Econometrica: Journal of the Econometric Society 424--438.
\item Gretton, A. (2015). Notes on mean embeddings and covariance operators. http:// www.gatsby.ucl.ac.uk/\textasciitilde{}gretton/coursefiles/lecture5\_covarianceOperator.pdf.
\item Gretton, A., Borgwardt, K. M., Rasch, M. J., Schölkopf, B. and Smola, A. (2012). A kernel two-sample test. Journal of Machine Learning Research 13:723--773.
\item Gretton, A., Bousquet, O., Smola, A. and Scholkopf, B. (2005). Measuring statistical dependence with Hilbert-Schmidt norms. Lecture Notes in Computer Science 63--78.
\item Guo, J., Levina, E., Michailidis, G. and Zhu, J. (2011). Joint estimation of multiple graphical models. Biometrika 98:1--15.
\item Hachiya, T., Furukawa, R., Shiwa, Y. et al. (2017). Genome-wide identification of inter-individually variable DNA methylation sites improves the efficacy of epigenetic association studies. npj Genomic Medicine 2:11.
\item Hansen, K. D., Wu, Z., Irizarry, R. A. and Leek, J. T. (2011). Sequencing technology does not eliminate biological variability. Nature Biotechnology 29:572--573.
\item Henderson, D. and Plaschko, P. (2006). Stochastic differential equations in science and engineering:(With CD-ROM). World Scientific.
\item Hernandez, K. M. (2015). Understanding the genetic architecture of complex traits using the function---Valued approach. New Phytologist 208:1--3.
\item Hoffmann, A., Ziller, M. and Spengler, D. (2016). The future is the past: Methylation QTLs in schizophrenia. Genes 7:104.
\item Hong, S., Chen, X., Jin, L. and Xiong, M. (2013). Canonical correlation analysis for RNA-seq co-expression networks. Nucleic Acids Research 41:e95.
\item Hosseini-Asl, E. (2016). Sparse feature learning for image analysis in segmentation, classification, and disease diagnosis. \url{https://doi.org/10.18297/etd/2456}.
\item Hoyer, P. O., Janzing, D., Mooij, J. M., Peters, J. and Schölkopf, B. (2009). Nonlinear causal discovery with additive noise models. Advances in Neural Information Processing Systems 689--696.
\item Hrdlickova, R., Toloue, M. and Tian, B. (2017). RNA-Seq methods for transcriptome analysis. Wiley Interdiscip Rev RNA 8.
\item Huan, T., Meng, Q., Saleh, M. A. et al. (2015). Integrative network analysis reveals molecular mechanisms of blood pressure regulation. Molecular Systems Biology 11: 799.
\item Huang, H.-C., Niu, Y. and Qin, L.-X. (2015). Differential expression analysis for RNA-seq: An overview of statistical methods and computational software. Cancer informatics 14:57--67.
\item Hyvärinen, A. (1998). New approximations of differential entropy for independent component analysis and projection pursuit. Advances in Neural Information Processing Systems 10:273--279.
\item Hyvärinen, A. and Smith, S. M. (2013). Pairwise likelihood ratios for estimation of non-Gaussian structural equation models. Journal of Machine Learning Research 14: 111--152.
\item Ideker, T. and Nussinov, R. (2017). Network approaches and applications in biology. PLoS Comput Biol 13:e1005771.
\item Jaakkola, T., Sontag, D., Globerson, A. and Meila, M. (2010). Learning Bayesian network structure using LP relaxations. Proceedings of the Thirteenth International Conference on Artificial Intelligence and Statistics 358--365.
\item Janzing, D., Hoyer, P. O. and Schölkopf, B. (2009). Telling cause from effect based on high-dimensional observations. In Proceedings of the 27th International Conference on Machine Learning pp. 479--486.
\item Janzing, D., Mooij, J., Zhang, K. et al. (2012). Information-geometric approach to inferring causal directions. Artificial Intelligence 182:1--31.
\item Janzing, D. and Scholkopf, B. (2010). Causal inference using the algorithmic Markov condition. IEEE Transactions on Information Theory 56:5168--5194.
\item Janzing, D., Steudel, B., Shajarisales, N. and Schölkopf, B. (2015). Justifying information-geometric causal inference. In Measures of complexity, Springer.
\item Javery, O., Shyn, P. and Mortele, K. (2013). FDG PET or PET/CT in patients with pancreatic cancer: When does it add to diagnostic CT or MRI? Clinical Imaging 37: 295--301.
\item Jefkine. (2016). Backpropagation in convolutional neural networks. \url{http://www.jefkine} .com/general/2016/09/05/backpropagation-in-convolutional-neural-networks.
\item Jiang, J., Lin, N., Guo, S., Chen, J. and Xiong, M. (2015). Multiple functional linear model for association analysis of RNA-seq with imaging. Quantitative Biology 3:90--102.
\item Jiang, P., Scarpa, J. R., Fitzpatrick, K. et al. (2015). A systems approach identifies networks and genes linking sleep and stress: Implications for neuropsychiatric disorders. Cell Reports 11:835--848.
\item Jin, B., Li, Y. and Robertson, K. D. (2011). DNA methylation: Superior or subordinate in the epigenetic hierarchy? Genes Cancer 2:607--617.
\item Jin, L.-p. and Dong, J. (2016). Ensemble deep learning for biomedical time series classification. Computational Intelligence and Neuroscience 2016:6212684.
\item Judge, G. G., Hill, R. C., Griffiths, W., Lutkepohl, H. and Lee, T. C. (1982). Introduction to the theory and practice of econometrics. John Wiley \& Sons, New York.
\item Köhler, S., Vasilevsky, N. A., Engelstad, M. et al. (2017). The human phenotype ontology in 2017. Nucleic Acids Research 45:D865--D876.
\item Kalisch, M. and Bühlmann, P. (2007). Estimating high-dimensional directed acyclic graphs with the PC-algorithm. Journal of Machine Learning Research 8:613--636.
\item Kalyagin, V. A., Koldanov, A. P., Koldanov, P. A. and Pardalos, P. M. (2017). Optimal statistical decision for Gaussian graphical model selection. arXiv:1701.02071.
\item Koivisto, M. and Sood, K. (2004). Exact Bayesian structure discovery in Bayesian networks. Journal of Machine Learning Research 5:549--573.
\item Kozlov, A. V. and Koller, D. (1997). Nonuniform dynamic discretization in hybrid networks. Proceedings of the Thirteenth Conference on Uncertainty in Artificial Intelligence 314--325.
\item Kpotufe, S., Sgouritsa, E., Janzing, D. and Schölkopf, B. (2014). Consistency of causal inference under the additive noise model. Proceedings of the 31st International Conference on Machine Learning (ICML-14) 478--486.
\item Krämer, N., Schäfer, J. and Boulesteix, A.-L. (2009). Regularized estimation of large-scale gene association networks using graphical Gaussian models. BMC Bioinformatics 10:384.
\item Kraskov, A., Stögbauer, H. and Grassberger, P. (2004). Estimating mutual information. Physical Review E 69:066138.
\item Kremling, A. and Saez-Rodriguez, J. (2007). Systems biology---An engineering perspective. Journal of Biotechnology 129:329--351.
\item Krizhevsky, A., Sutskever, I. and Hinton, G. E. (2012). Imagenet classification with deep convolutional neural networks. Advances in Neural Information Processing Systems 1097--1105.
\item Krueger, F., Kreck, B., Franke, A. and Andrews, S. R. (2012). DNA methylome analysis using short bisulfite sequencing data. Nature Methods 9:145--151.
\item Kulig, P., Pach, R. and Kulig, J. (2014). Role of abdominal ultrasonography in clinical staging of pancreatic carcinoma: A tertiary center experience. Pol Arch Med Wewn 124:225--232.
\item Kurdyukov, S. and Bullock, M. (2016). DNA Methylation Analysis: Choosing the Right Method. Biology (Basel) 5:3.
\item Kwak, I.-Y., Moore, C. R., Spalding, E. P. and Broman, K. W. (2014). A simple regression-based method to map quantitative trait loci underlying function-valued phenotypes. Genetics 197:1409--1416.
\item Kwak, I.-Y., Moore, C. R., Spalding, E. P. and Broman, K. W. (2016). Mapping quantitative trait loci underlying function-valued traits using functional principal component analysis and multi-trait mapping. G3: Genes, Genomes, Genetics 6:79--86.
\item Laird, P. W. (2010). Principles and challenges of genome-wide DNA methylation analysis. Nature Reviews Genetics 11:191--203.
\item Lappalainen, T. and Greally, J. M. (2017). Associating cellular epigenetic models with human phenotypes. Nat Rev Genet 18:441--451.
\item Lappalainen, T., Sammeth, M., Friedländer, M. R. et al. (2013). Transcriptome and genome sequencing uncovers functional variation in humans. NATURE 501:506--511.
\item LeCun, Y. (1989). Generalization and network design strategies. Connectionism in perspective 143--155.
\item LeCun, Y., Bengio, Y. and Hinton, G. (2015). Deep learning. NATURE 521:436--444.
\item Lee, D. D. and Seung, H. S. (1999). Learning the parts of objects by non-negative matrix factorization. NATURE 401:788--791.
\item Lee, D. D. and Seung, H. S. (2001). Algorithms for non-negative matrix factorization. Advances in Neural Information Processing Systems 556--562.
\item Lee, D. Y. (2015). Nonlinear functional regression model for sequencing-based association and gene-gene interaction analysis of physiological traits and their applications to sleep apnea. Thesis, The University of Texas School of Public Health.
\item Lelli, K. M., Slattery, M. and Mann, R. S. (2012). Disentangling the many layers of eukaryotic transcriptional regulation. Annual Review of Genetics 46:43--68.
\item Lemeire, J. and Janzing, D. (2013). Replacing causal faithfulness with algorithmic independence of conditionals. Minds and Machines 23:227--249.
\item Li, B. and Dewey, C. N. (2011). RSEM: Accurate transcript quantification from RNA-Seq data with or without a reference genome. BMC Bioinformatics 12:323.
\item Li, B., Ruotti, V., Stewart, R. M., Thomson, J. A. and Dewey, C. N. (2009). RNA-Seq gene expression estimation with read mapping uncertainty. Bioinformatics 26: 493--500.
\item Li, D., Xie, Z., Pape, M. L. and Dye, T. (2015). An evaluation of statistical methods for DNA methylation microarray data analysis. BMC Bioinformatics 16:217.
\item Li, J., Witten, D. M., Johnstone, I. M. and Tibshirani, R. (2012). Normalization, testing, and false discovery rate estimation for RNA-sequencing data. Biostatistics 13: 523--538.
\item Li, L. and Xiong, M. (2015). Dynamic model for RNA-seq data analysis. BioMed Research International 2015:916352.
\item Li, P., Piao, Y., Shon, H. S. and Ryu, K. H. (2015). Comparing the normalization methods for the differential analysis of Illumina high-throughput RNA-Seq data. BMC Bioinformatics 16:347.
\item Li, X., Dunn, J., Salins, D. et al. (2017). Digital health: Tracking physiomes and activity using wearable biosensors reveals useful health-related information. PLoS Biology 15:e2001402.
\item Li, Y., Nan, B. and Zhu, J. (2015). Multivariate sparse group lasso for the multivariate multiple linear regression with an arbitrary group structure. Biometrics 71:354--363.
\item Lienart, T. (2015). RKHS Embeddings. \url{https://www.stats.ox.ac.uk/~lienart/gml15} \_rkhsembeddings.html.
\item Lin, D., Calhoun, V. D. and Wang, Y.-P. (2014). Correspondence between fMRI and SNP data by group sparse canonical correlation analysis. Med Image Anal 18:891-- 902.
\item Lin, G., Shen, C., van den Hengel, A. and Reid, I. (2016). Efficient piecewise training of deep structured models for semantic segmentation. Proceedings of the IEEE Conference on Computer Vision and Pattern Recognition 3194--3203.
\item Lin, N., Jiang, J., Guo, S. and Xiong, M. (2015). Functional principal component analysis and randomized sparse clustering algorithm for medical image analysis. PLoS One 10:e0132945.
\item Litjens, G., Kooi, T., Bejnordi, B. E. et al. (2017). A survey on deep learning in medical image analysis. Med Image Anal 42:60--88.
\item Liu, F. and Chan, L. (2016). Causal inference on discrete data via estimating distance correlations. Neural Computation 28:801--814.
\item Liu, F., Zhang, S.-W., Guo, W.-F., Wei, Z.-G. and Chen, L. (2016). Inference of gene regulatory network based on local bayesian networks. PLoS Comput Biol 12: e1005024.
\item Liu, G., Lin, Z., Yan, S., Sun, J., Yu, Y. and Ma, Y. (2013). Robust recovery of subspace structures by low-rank representation. IEEE Transactions on Pattern Analysis and Machine Intelligence 35:171--184.
\item Liu, X., Shi, X., Chen, C. and Zhang, L. (2015). Improving RNA-Seq expression estimation by modeling isoform-and exon-specific read sequencing rate. BMC Bioinformatics 16:332.
\item Long, J., Shelhamer, E. and Darrell, T. (2015). Fully convolutional networks for semantic segmentation. Proceedings of the IEEE Conference on Computer Vision and Pattern Recognition 3431--3440.
\item Louizos, C., Shalit, U., Mooij, J., Sontag, D., Zemel, R. and Welling, M. (2017). Causal effect inference with deep latent-variable models. arXiv:1705.08821.
\item Luo, L., Zhu, Y. and Xiong, M. (2013). Smoothed functional principal component analysis for testing association of the entire allelic spectrum of genetic variation. European Journal of Human Genetics 21:217--224.
\item Ma, C.-X., Casella, G. and Wu, R. (2002). Functional mapping of quantitative trait loci underlying the character process: A theoretical framework. Genetics 161:1751-- 1762.
\item Magliacane, S., van Ommen, T., Claassen, T., Bongers, S., Versteeg, P. and Mooij, J. M. (2017). Causal transfer learning. arXiv:1707.06422.
\item Mairal, J. and Yu, B. (2013). Supervised feature selection in graphs with path coding penalties and network flows. The Journal of Machine Learning Research 14:2449-- 2485.
\item Marblestone, A. H., Wayne, G. and Kording, K. P. (2016). Toward an integration of deep learning and neuroscience. Frontiers in Computational Neuroscience 10:94.
\item Mazumder, R. and Hastie, T. (2012). The graphical lasso: New insights and alternatives. Electronic Journal of Statistics 6:2125--2149.
\item McCarthy, D. J., Chen, Y. and Smyth, G. K. (2012). Differential expression analysis of multifactor RNA-Seq experiments with respect to biological variation. Nucleic Acids Research 40:4288--4297.
\item McRae, A., Marioni, R. E., Shah, S. et al. (2017). Identification of 55,000 Replicated DNA Methylation QTL. bioRxiv: 166710.
\item Medland, S. E., Jahanshad, N., Neale, B. M. and Thompson, P. M. (2014). Whole-genome analyses of whole-brain data: Working within an expanded search space. Nature Neuroscience 17:791--800.
\item Meier, L., Van de Geer, S. and Bühlmann, P. (2009). High-dimensional additive modeling. The Annals of Statistics 37:3779--3821.
\item Menéndez, P., Kourmpetis, Y. A., ter Braak, C. J. and van Eeuwijk, F. A. (2010). Gene regulatory networks from multifactorial perturbations using Graphical Lasso: Application to the DREAM4 challenge. PLoS One 5:e14147.
\item Mi, X., Eskridge, K., Wang, D. et al. (2010). Regression-based multi-trait QTL mapping using a structural equation model. Statistical Applications in Genetics and Molecular Biology 9:Article38.
\item Moeskops, P., de Bresser, J., Kuijf, H. J. et al. (2018). Evaluation of a deep learning approach for the segmentation of brain tissues and white matter hyperintensities of presumed vascular origin in MRI. NeuroImage: Clinical 17:251--262.
\item Mohammadi, S., Zuckerman, N., Goldsmith, A. and Grama, A. (2017). A Critical Survey of Deconvolution Methods for Separating Cell Types in Complex Tissues. Proceedings of the IEEE 105:340--366.
\item Mohan, K., London, P., Fazel, M., Witten, D. and Lee, S.-I. (2014). Node-based learning of multiple gaussian graphical models. The Journal of Machine Learning Research 15:445--488.
\item Moneta, A., Entner, D., Hoyer, P. O. and Coad, A. (2013). Causal inference by independent component analysis: Theory and applications. Oxford Bulletin of Economics and Statistics 75:705--730.
\item Mooij, J. and Janzing, D. (2010). Distinguishing between cause and effect. Causality: Objectives and Assessment 147--156.
\item Mooij, J. M., Peters, J., Janzing, D., Zscheischler, J. and Schölkopf, B. (2016). Distinguishing cause from effect using observational data: Methods and benchmarks. The Journal of Machine Learning Research 17:1103--1204.
\item Nariai, N., Hirose, O., Kojima, K. and Nagasaki, M. (2013). TIGAR: Transcript isoform abundance estimation method with gapped alignment of RNA-Seq data by variational Bayesian inference. Bioinformatics 29:2292--2299.
\item Nariai, N., Kojima, K., Mimori, T., Kawai, Y. and Nagasaki, M. (2016). A Bayesian approach for estimating allele-specific expression from RNA-Seq data with diploid genomes. BMC Genomics 17:Suppl 1:2.
\item Nelms, B. D., Waldron, L., Barrera, L. A. et al. (2016). CellMapper: Rapid and accurate inference of gene expression in difficult-to-isolate cell types. Genome Biol 17:201.
\item Nesterov, Y. and Nemirovskii, A. (1994). Interior-point polynomial algorithms in convex programming. SIAM.
\item Noh, H. and Gunawan, R. (2016). Inferring gene targets of drugs and chemical compounds from gene expression profiles. Bioinformatics 32:2120--2127.
\item Noh, H., Hong, S. and Han, B. (2015). Learning deconvolution network for semantic segmentation. Proceedings of the IEEE International Conference on Computer Vision 1520--1528.
\item Nowozin, S. and Lampert, C. H. (2011). Structured learning and prediction in computer vision. Foundations and Trends® in Computer Graphics and Vision 6:185--365.
\item Nowzohour, C. and Bühlmann, P. (2016). Score-based causal learning in additive noise models. Statistics 50:471--485.
\item Ogata, K. (1998). System dynamics. 3rd ed., Prentice Hall, New Jersey.
\item Orho-Melander, M. (2015). Genetics of coronary heart disease: Towards causal mechanisms, novel drug targets and more personalized prevention. Journal of Internal Medicine 278:433--446.
\item Parikh, N. and Boyd, S. (2014). Proximal algorithms. Foundations and Trends® in Optimization 1:127--239.
\item Pearl, J. (2009). Causality: Models, reasoning, and inference. 2nd ed., Cambridge University Press, New York.
\item Pearl, J. (2013). The mathematics of causal inference. Proceedings of the Joint Statistical Meetings Conference.
\item Pergola, G., Selvaggi, P., Trizio, S., Bertolino, A. and Blasi, G. (2015). The role of the thalamus in schizophrenia from a neuroimaging perspective. Neuroscience \& Biobehavioral Reviews 54:57--75.
\item Peters, J. and Bühlmann, P. (2013). Identifiability of Gaussian structural equation models with equal error variances. Biometrika 101:219--228.
\item Peters, J., Janzing, D. and Schölkopf, B. (2017). Elements of causal inference: Foundations and learning algorithms. MIT Press, Boston.
\item Peters, J., Janzing, D. and Schölkopf, B. (2011). Causal inference on discrete data using additive noise models. IEEE Transactions on Pattern Analysis and Machine Intelligence 33:2436--2450.
\item Peters, J., Mooij, J., Janzing, D. and Schölkopf, B. (2012). Identifiability of causal graphs using functional models. arXiv:1202.3757.
\item Peters, J. M. (2012). Restricted structural equation models for causal inference, Ph.D. Thesis, ETH Zurich.
\item Pettit, J. B., Tomer, R., Achim, K., Richardson, S., Azizi, L. and Marioni, J. (2014). Identifying cell types from spatially referenced single-cell expression datasets. PLoS Comput Biol 10:e1003824.
\item Pidsley, R., Zotenko, E., Peters, T. J. et al. (2016). Critical evaluation of the Illumina MethylationEPIC BeadChip microarray for whole-genome DNA methylation profiling. Genome Biol 17:208.
\item Poyton, A., Varziri, M. S., McAuley, K. B., McLellan, P. and Ramsay, J. O. (2006). Parameter estimation in continuous-time dynamic models using principal differential analysis. Computers \& Chemical Engineering 30:698--708.
\item Rakyan, V. K., Down, T. A., Balding, D. J. and Beck, S. (2011). Epigenome-wide association studies for common human diseases. Nat Rev Genet 12:529--541.
\item Ramsay, J. and Silverman, B. W. (2005). Functional data analysis. 2nd ed., Springer-Verlag, New York.
\item Rapin, J., Bobin, J., Larue, A. and Starck, J. L. (2012). Robust non-negative matrix factorization for multispectral data with sparse prior. Proceedings of ADA7.
\item Ravì, D., Wong, C., Lo, B. and Yang, G.-Z. (2017). A deep learning approach to on-node sensor data analytics for mobile or wearable devices. IEEE Journal of Biomedical and Health Informatics 21:56--64.
\item Ravier, P., Leclerc, F., Dumez-Viou, C. and Lamarque, G. (2007). Redefining performance evaluation tools for real-time QRS complex classification systems. IEEE Transactions on Biomedical Engineering 54:1706--1710.
\item Richfield, O., Alam, M. A., Calhoun, V. and Wang, Y.-P. (2016). Learning Schizophrenia Imaging Genetics Data Via Multiple Kernel Canonical Correlation Analysis. Bioinformatics and Biomedicine (BIBM), 2016 IEEE International Conference on 507--511.
\item Robinson, M. D. and Oshlack, A. (2010). A scaling normalization method for differential expression analysis of RNA-seq data. Genome Biol 11:R25.
\item Ronneberger, O., Fischer, P. and Brox, T. (2015). U-net: Convolutional networks for biomedical image segmentation. International Conference on Medical Image Computing and Computer-Assisted Intervention 234--241.
\item Rosa, G. J., Valente, B. D., de los Campos, G., Wu, X.-L., Gianola, D. and Silva, M. A. (2011). Inferring causal phenotype networks using structural equation models. Genetics Selection Evolution 43:6.
\item Ross, S. M. (2014). Introduction to probability models. 11th ed., Academic Press, San Diego.
\item Sagan, H. (2012). Introduction to the calculus of variations. Courier Corporation.
\item Sandler, R. (2010). Nonnegative matrix factorization for segmentation analysis. Technion-Israel Institute of Technology, Faculty of Computer Science.
\item Sason, I. and Verdú, S. (2016). f -Divergence Inequalities. IEEE Transactions on Information Theory 62:5973--6006.
\item Sathyanarayana, A., Joty, S., Fernandez-Luque, L. et al. (2016). Sleep quality prediction from wearable data using deep learning. JMIR mHealth and uHealth 4:e125.
\item Sauwen, N., Acou, M., Sima, D. M. et al. (2017). Semi-automated brain tumor segmentation on multi-parametric MRI using regularized non-negative matrix factorization. BMC Medical Imaging 17:29.
\item Scanagatta, M., De Campos, C. P. and Zaffalon, M. (2014). Min-BDeu and max-BDeu scores for learning Bayesian networks. European Workshop on Probabilistic Graphical Models 426--441.
\item Scutari, M. (2017). Package 'bnlearn'. \url{http://www.bnlearn.com}.
\item Shen-Orr, S. S. and Gaujoux, R. (2013). Computational deconvolution: Extracting cell type-specific information from heterogeneous samples. Current Opinion in Immunology 25:571--578.
\item Shim, H. and Stephens, M. (2015). Wavelet-based genetic association analysis of functional phenotypes arising from high-throughput sequencing assays. The Annals of Applied Statistics 9:665--686.
\item Shimizu, S., Hoyer, P. O., Hyvärinen, A. and Kerminen, A. (2006). A linear non-Gaussian acyclic model for causal discovery. Journal of Machine Learning Research 7:2003--2030.
\item Shimizu, S., Inazumi, T., Sogawa, Y. et al. (2011). DirectLiNGAM: A direct method for learning a linear non-Gaussian structural equation model. Journal of Machine Learning Research 12:1225--1248.
\item Shrout, P., Keyes, K. and Ornstein, K. (2011). Causality and psychopathology: Finding the determinants of disorders and their cures. Oxford University Press.
\item Sims, C. A. (1972). Money, income, and causality. The American Economic Review 62: 540--552.
\item Skelly, D. A., Johansson, M., Madeoy, J., Wakefield, J. and Akey, J. M. (2011). A powerful and flexible statistical framework for testing hypotheses of allele-specific gene expression from RNA-seq data. Genome Research 21:1728--1737.
\item Spirtes, P., Meek, C., Richardson, T. and Meek, C. (1999). An algorithm for causal inference in the presence of latent variables and selection bias.
\item Statnikov, A., Henaff, M., Lytkin, N. I. and Aliferis, C. F. (2012). New methods for separating causes from effects in genomics data. BMC Genomics 13:S22.
\item Stegle, O., Teichmann, S. A. and Marioni, J. C. (2015). Computational and analytical challenges in single-cell transcriptomics. Nature Reviews Genetics 16:133--145.
\item Stein, J. L., Hua, X., Lee, S. et al. (2010). Voxelwise genome-wide association study (vGWAS). neuroimage 53:1160--1174.
\item Stephens, M. (2013). A unified framework for association analysis with multiple related phenotypes. PLoS One 8:e65245.
\item Su, J., Yi, D., Liu, C., Guo, L. and Chen, W.-H. (2017). Dimension reduction aided hyperspectral image classification with a small-sized training dataset: Experimental comparisons. Sensors 17:2726.
\item Sun, A., Venkatesh, A. and Hall, D. A. (2016). A multi-technique reconfigurable electrochemical biosensor: Enabling personal health monitoring in mobile devices. IEEE Transactions on Biomedical Circuits and Systems 10:945--954.
\item Sun, K., Gonçalves, J. P., Larminie, C. and Pržulj, N. (2014). Predicting disease associations via biological network analysis. BMC Bioinformatics 15:304.
\item Sun, Z. and Zhu, Y. (2012). Systematic comparison of RNA-Seq normalization methods using measurement error models. Bioinformatics 28:2584--2591.
\item Sutton, C. A. and McCallum, A. (2005). Piecewise training for undirected models. Proceeding UAI'05 Proceedings of the Twenty-First Conference on Uncertainty in Artificial Intelligence. Pp. 568--575. Edinburgh, Scotland --- July 26--29, 2005.
\item Sutton, C. and McCallum, A. (2011). An introduction to conditional random fields. Foundations and Trends® in Machine Learning 4:267--373.
\item Székely, G. J. and Rizzo, M. L. (2009). Brownian distance covariance. The Annals of Applied Statistics 3:1236--1265.
\item Székely, G. J., Rizzo, M. L. and Bakirov, N. K. (2007). Measuring and testing dependence by correlation of distances. The Annals of Statistics 35:2769--2794.
\item Takahashi, T., Zhou, S.-Y., Nakamura, K. et al. (2011). A follow-up MRI study of the fusiform gyrus and middle and inferior temporal gyri in schizophrenia spectrum. Progress in Neuro-Psychopharmacology and Biological Psychiatry 35:1957--1964.
\item Tian, Y., Morris, T. J., Webster, A. P. et al. (2017). ChAMP: Updated methylation analysis pipeline for Illumina BeadChips. Bioinformatics 33:3982--3984.
\item Tibshirani, R., Saunders, M., Rosset, S., Zhu, J. and Knight, K. (2005). Sparsity and smoothness via the fused lasso. Journal of the Royal Statistical Society: Series B (Statistical Methodology) 67:91--108.
\item Trapnell, C., Hendrickson, D. G., Sauvageau, M., Goff, L., Rinn, J. L. and Pachter, L. (2013). Differential analysis of gene regulation at transcript resolution with RNA-seq. Nature Biotechnology 31:46--53.
\item Udell, M., Horn, C., Zadeh, R. and Boyd, S. (2016). Generalized low rank models. Foundations and Trends® in Machine Learning 9:1--118.
\item Valente, B. D. and de Magalhaes Rosa, G. J. (2013). Mixed effects structural equation models and phenotypic causal networks. Genome-Wide Association Studies and Genomic Prediction 449--464.
\item Varvarigou, V., Dahabreh, I. J., Malhotra, A. and Kales, S. N. (2011). A review of genetic association studies of obstructive sleep apnea: Field synopsis and meta-analysis. Sleep 34:1461--1468.
\item Vavasis, S. A. (2009). On the complexity of nonnegative matrix factorization. SIAM Journal on Optimization 20:1364--1377.
\item Verma, T. and Pearl, J. (1991). Equivalence and synthesis of causal models. In Proceedings of the 6th Annual Conference on Uncertainty in Artificial Intelligence (UAI) 255--270.
\item Vincent, M., Mundbjerg, K., Skou Pedersen, J. et al. (2017). epiG: Statistical inference and profiling of DNA methylation from whole-genome bisulfite sequencing data. Genome Biol 18:38.
\item Wang, B., Huang, L., Zhu, Y., Kundaje, A., Batzoglou, S. and Goldenberg, A. (2017). Vicus: Exploiting local structures to improve network-based analysis of biological data. PLoS Comput Biol 13:e1005621.
\item Wang, H., Nie, F., Huang, H. et al. (2011). Identifying quantitative trait loci via group-sparse multitask regression and feature selection: An imaging genetics study of the ADNI cohort. Bioinformatics 28:229--237.
\item Wang, J., Ding, H., Azamian, F. et al. (2017). Detecting cardiovascular disease from mammograms with deep learning. IEEE Transactions on Medical Imaging 36:1172-- 1181.
\item Wang, K., Zhao, Y., Xiong, Q. et al. (2016). Research on healthy anomaly detection model based on deep learning from multiple time-series physiological signals. Scientific Programming 2016:5642856.
\item Wang, P. (2016). Causal genetic network analysis for multiple phenotypes from complex diseases. Thesis, Fudan University.
\item Wang, P., Rahman, M., Jin, L. and Xiong, M. (2016). A new statistical framework for genetic pleiotropic analysis of high dimensional phenotype data. BMC genomics 17:881.
\item Wang, T. and Zhu, L. (2013). Sparse sufficient dimension reduction using optimal scoring. Computational Statistics \& Data Analysis 57:223--232.
\item Wang, Y. (2011). Smoothing splines: Methods and applications. CRC Press, New York.
\item Wang, Z., Fang, H., Tang, N. L.-S. and Deng, M. (2017). VCNet: Vector-based gene co-expression network construction and its application to RNA-seq data. Bioinformatics 33:2173--2181.
\item Wei, B., Sun, X., Ren, X. and Xu, J. (2017). Minimal effort back propagation for convolutional neural networks. arXiv:1709.05804.
\item Wolfe, E., Spekkens, R. W. and Fritz, T. (2016). The inflation technique for causal inference with latent variables. arXiv:1609.00672.
\item Wright, S. (1921). Correlation and causation. Journal of Agricultural Research 20:557--585.
\item Wu, J., Pan, S., Zhu, X. and Cai, Z. (2015). Boosting for multi-graph classification. IEEE Transactions on Cybernetics 45:416--429.
\item Wu, M. C., Lee, S., Cai, T., Li, Y., Boehnke, M. and Lin, X. (2011). Rare-variant association testing for sequencing data with the sequence kernel association test. Am J Hum Genet 89:82--93.
\item Xie, Y., Zhang, Z., Sapkota, M. and Yang, L. (2016). Spatial clockwork recurrent neural network for muscle perimysium segmentation. International Conference on Medical Image Computing and Computer-Assisted Intervention 185--193.
\item Xiong, H., Brown, J. B., Boley, N., Bickel, P. J. and Huang, H. (2014). DE-FPCA: Testing gene differential expression and exon usage through functional principal component analysis. In Statistical analysis of next generation sequencing data: Springer.
\item Xiong, M. M. (2018). Big data in omics and imaging: Association analysis. Chapman and Hall/CRC.
\item Xu, K., Jin, L. and Xiong, M. (2017). Functional regression method for whole genome eQTL epistasis analysis with sequencing data. BMC Genomics 18:385.
\item Yang, J., Wu, R. and Casella, G. (2009). Nonparametric functional mapping of quantitative trait loci. Biometrics 65:30--39.
\item Yong, W. S., Hsu, F. M. and Chen, P. Y. (2016). Profiling genome-wide DNA methylation. Epigenetics Chromatin 9:26.
\item Yu, Q., Erhardt, E. B., Sui, J. et al. (2015). Assessing dynamic brain graphs of time-varying connectivity in fMRI data: Application to healthy controls and patients with schizophrenia. Neuroimage 107:345--355.
\item Yuan, M. and Lin, Y. (2006). Model selection and estimation in regression with grouped variables. Journal of the Royal Statistical Society: Series B (Statistical Methodology) 68:49--67.
\item Zeiler, M. D., Krishnan, D., Taylor, G. W. and Fergus, R. (2010). Deconvolutional networks. Computer Vision and Pattern Recognition (CVPR), 2010 IEEE Conference on 2528--2535.
\item Zeiler, M. D., Taylor, G. W. and Fergus, R. (2011). Adaptive deconvolutional networks for mid and high level feature learning. Computer Vision (ICCV), 2011 IEEE International Conference on 2018--2025.
\item Zhang, B., Gaiteri, C., Bodea, L.-G. et al. (2013). Integrated systems approach identifies genetic nodes and networks in late-onset Alzheimer's disease. Cell 153: 707--720.
\item Zhang, K., Wang, Z., Zhang, J. and Schölkopf, B. (2016). On estimation of functional causal models: General results and application to the post-nonlinear causal model. ACM Transactions on Intelligent Systems and Technology (TIST) 7:1:22.
\item Zhang, Q., Filippi, S., Gretton, A. and Sejdinovic, D. (2017). Large-scale kernel methods for independence testing. Statistics and Computing pp. 1--18.
\item Zhang, Z. (2016). Derivation of Backpropagation in Convolutional Neural Network (CNN). \url{http://web.eecs.utk.edu/~zzhang61}.
\item Zhang, Z. H., Jhaveri, D. J., Marshall, V. M. et al. (2014). A comparative study of techniques for differential expression analysis on RNA-Seq data. PLoS One 9: e103207.
\item Zhao, S., Zhang, Y., Gordon, W. et al. (2015). Comparison of stranded and non-stranded RNA-seq transcriptome profiling and investigation of gene overlap. BMC Genomics 16:675.
\item Zheng, Y., Liu, Q., Chen, E., Ge, Y. and Zhao, J. L. (2016). Exploiting multi-channels deep convolutional neural networks for multivariate time series classification. Frontiers of Computer Science 10:96--112.
\item Zhou, S., Rütimann, P., Xu, M. and Bühlmann, P. (2011). High-dimensional covariance estimation based on Gaussian graphical models. Journal of Machine Learning Research 12:2975--3026.
\item Zhou, X. and Stephens, M. (2014). Efficient multivariate linear mixed model algorithms for genome-wide association studies. Nature Methods 11:407--409.
\item Zhou, Y., Arpit, D., Nwogu, I. and Govindaraju, V. (2014). Is Joint Training Better for Deep Auto-Encoders? arXiv:1405.1380.
\item Zhu, H., Rao, R. S. P., Zeng, T. and Chen, L. (2012). Reconstructing dynamic gene regulatory networks from sample-based transcriptional data. Nucleic Acids Research 40:10657--10667.
\item Zscheischler, J., Janzing, D. and Zhang, K. (2012). Testing whether linear equations are causal: A free probability theory approach. arXiv:1202.3779.
\item Zyprych-Walczak, J., Szabelska, A., Handschuh, L. et al. (2015). The impact of normalization methods on RNA-Seq data analysis. BioMed Research International 2015:621690.
\end{list}
\endgroup
